\documentclass[10pt]{article}
\usepackage{icmlko}

\icmlmeta{IP 파운데이션 월드모델}{269}{도서는 시계였다}{네 노트가 따로 짚은 도서 이상이 하나의 사실이었다 --- 유보 안에서만 라벨이 8.7배 쌓인다}

\begin{document}

\twocolumn[
\icmlheader

\begin{abstract}
도서가 네 번 짚혔다. 노트 252는 여지 $+$0.177 로 아홉 중 1위라 했고, 노트
253은 \textbf{제 학습 자료를 덜어 내면 오히려 좋아지는} 유일한 도메인임을
쟀고(80건 $\to$ 40건에서 $\rho$ 0.321 $\to$ 0.360), 노트 254는 여덟 도메인에
같은 검사를 돌려 도서만 그렇다고 확인했고, 노트 267은 \textbf{``$-$시작일''
한 줄이 축 스물셋보다 낫다}는 것을 봤다(0.5615 대 0.3480). 네 번 다 원인을
``베스트셀러 스크랩이라 그런가 보다''로 넘겼다. 이번에 라벨을 직접 열었다.
\textbf{유보 163건 \emph{안에서만} 반년 코호트별 sales point 중앙값이
39,209 $\to$ 28,991 $\to$ 18,070 $\to$ 4,485 다} --- 1년 반 사이 \textbf{8.7배}.
\texttt{review\_count} 도 출간 시점과 $+$0.637 로 같은 무늬다. \textbf{도서
라벨은 ``얼마나 팔렸나''가 아니라 ``언제 나왔나''를 재고 있다.} 네 이상이
하나의 사실에서 나온다 --- 여지가 큰 것도(모형이 시계를 못 따라가서),
제 학습이 해로운 것도(1999--2025 학습은 다른 시계다), 달력이 이기는 것도.
\textbf{고칠 방법은 둘인데 둘 다 지금 못 한다.} 게임처럼 고정 창 라벨로
바꾸려면(노트 210의 \texttt{y\_w30}) 새로 긁어야 하고, 코호트 안 순위로
라벨을 다시 정의하는 것은 \textbf{과제를 바꾸는 결정}이라 점수로 정하면 안
된다(노트 224). \textbf{그래서 안 고치고, 대신 도서 수를 읽는 법을 적어
둔다} --- 도서 $\rho$ 0.35 는 능력이 아니라 대부분 시계다.
\end{abstract}
]

\section{네 번 짚고 네 번 다 넘겼다}

\begin{center}\small
\begin{tabular}{lll}
\hline
노트 & 무엇 & 그때 쓴 원인 \\
\hline
252 & 여지 $+$0.177 (1위) & ``학습 80 대 유보 163'' \\
253 & 학습을 덜면 좋아짐 & ``베스트셀러 스크랩'' \\
254 & 아홉 중 도서만 & ``다른 모집단'' \\
267 & 달력이 모형을 이김 & ``날짜가 곧 라벨'' \\
\hline
\end{tabular}
\end{center}

넷 다 그럴듯했고 넷 다 \emph{라벨을 직접 안 열어 보고} 쓴 말이다. 이번에
열었다.

\section{유보 안에서만 8.7배}

\begin{figure}[t]
\centering
\includegraphics[width=\columnwidth]{figs/book.pdf}
\caption{왼쪽 --- 유보 163건을 반년 코호트로 나눈 sales point 중앙값.
오른쪽 --- 도서가 짚힌 네 자리.}
\end{figure}

\begin{center}\small
\begin{tabular}{lrrr}
\hline
출간 & 건수 & sales point 중앙 & 안에서 $-$출간일 \\
\hline
2025 H1 & 12 & \textbf{39,209} & $-$0.173 \\
2025 H2 & 18 & 28,991 & $-$0.031 \\
2026 H1 & 79 & 18,070 & $+$0.206 \\
2026 H2$+$ & 54 & \textbf{4,485} & $+$0.189 \\
\hline
\end{tabular}
\end{center}

\textbf{1년 반 사이 8.7배다.} 그리고 코호트 \emph{안에서는} $-$출간일의
상관이 $-$0.17 $\sim$ $+$0.21 로 흩어진다 --- 즉 코호트 사이의 차이가 거의
전부다.

같은 무늬가 다른 필드에도 있다.

\begin{center}\small
\begin{tabular}{lr}
\hline
 & spearman \\
\hline
\texttt{review\_count} 대 sales point & $+$0.721 \\
\texttt{review\_count} 대 $-$출간일 & $+$0.637 \\
\hline
\end{tabular}
\end{center}

리뷰 수도 같이 쌓인다. \textbf{도서 레코드에서 ``인기''로 보이는 것이 전부
시계다.}

\section{네 이상이 한 사실이다}

이제 넷이 한 줄로 꿰인다.

\begin{itemize}
\item \textbf{여지가 1위}(252) --- 유보 안 CV 모형은 유보 라벨을 보고 배우므로
시계를 그대로 학습한다. 우리 모형은 못 하니 격차가 커 보인다. \emph{도달할
수 있는 여지가 아니라 시계의 크기였다.}
\item \textbf{제 학습이 해롭다}(253 · 254) --- 학습 80건은 1999--2025 에
흩어져 있고 그 시계는 유보의 시계와 기울기가 다르다. 옛 책은 20년 쌓였고
2026년 책은 두 달 쌓였다. \emph{다른 시계를 배우면 해롭다.}
\item \textbf{달력이 모형을 이긴다}(267) --- 라벨이 시계면 시계가 제일 좋은
예측이다. 당연한 결과였다.
\end{itemize}

\section{고칠 방법 둘, 둘 다 못 한다}

\paragraph{고정 창 라벨} 노트 210이 게임에서 한 것이다 --- 누적 리뷰
(\texttt{y\_raw})를 출시 30일 창(\texttt{y\_w30})으로 바꾸니 시작연도 상관이
$-$0.648에서 $-$0.071 로 사라지고 게임 $\rho$ 가 올랐다. 도서에는 그런
필드가 \textbf{없다}(노트 267에서 확인했다 --- 아홉 중 게임에만 있다).
새로 긁어야 하는데 알라딘 판매지수는 시계열 이력을 안 준다.

\paragraph{코호트 안 순위로 라벨 재정의} 값싸고 지금 할 수 있다. 그런데
그것은 \textbf{과제를 바꾸는 결정}이다 --- ``얼마나 팔릴까''를 ``같은 달에
나온 책들 중 몇 등일까''로 바꾼다. 노트 224가 평가 점수 축을 물리며 적은
규칙이 그대로 걸린다: \textbf{과제를 바꾸는 결정은 점수로 정하지 않는다.}
게다가 노트 268이 코호트 눈금으로 줄을 세워도 순위가 안 바뀐다는 것을
쟀으니, 바꿔서 얻을 것도 분명하지 않다.

\textbf{그래서 안 고친다.} 대신 읽는 법을 적어 둔다.

\section{도서 수를 읽는 법}

\begin{quote}
\textbf{도서 $\rho$ 0.35 는 능력이 아니라 대부분 시계다.} 도서를 올리는
변경은 ``시계를 더 잘 맞히게 됐다''일 수 있으므로, 도서만 오르는 결과는
믿지 않는다. 도서 몫이 6.2\%이므로 판에 미치는 영향은 $\pm$0.01 안쪽이다.
\end{quote}

이것은 노트 209가 판 전체에 적어 둔 것(``판 $\rho$ 의 0.25 는 언제 나왔나
다'')의 도메인 판이고, 노트 267이 판 수준에서는 그 걱정이 작다는 것을
쟀다(모형은 코호트 통제에 1\%만 준다). \textbf{판 수준에서 작은 것이
도서에서는 지배적이다.}

\paragraph{남은 것}
\begin{center}\small
\begin{tabular}{ll}
\hline
갈래 & 왜 \\
\hline
도서 재수집 & 고정 창이 있어야 고쳐진다 \\
같은 검사 아홉 곳 & 코호트별 라벨 중앙값 비를 다 재면 \\
팝업 · 펀딩 & 코호트 조각이 하나라 이 검사를 못 한다 \\
\texttt{SOURCE} & 34/45 쌍 \\
\hline
\end{tabular}
\end{center}

\paragraph{규약} 도메인이 두 번 넘게 이상하게 나오면 \textbf{라벨을 직접
열어 본다.} 도서는 네 번 짚히는 동안 축과 학습 표본만 봤고, 원인은 라벨
한 줄에 있었다. ``덜어 내 본다''(253) · ``잡음을 먹여 본다''(259) ·
``떼어 본다''(261) · ``코호트 안에서 재 본다''(267)에 이어 다섯 번째 ---
\textbf{코호트별 라벨 중앙값을 적어 본다.}

\end{document}
